\documentclass[11pt,a4paper]{article}
\usepackage[utf8]{inputenc}
\usepackage[T1]{fontenc}
\usepackage[spanish]{babel}
\usepackage{geometry}
\usepackage{verbatim}
\geometry{margin=2.5cm}

\title{Analizador léxico y sintáctico de C-TDS \\ \large Diseño e interfaz del AST}
\author{Agustin Alieni, Fran Natale y Julian Varea}
\date{}

\begin{document}
\maketitle

\section{Introducción}

En esta etapa construimos el analizador léxico (\texttt{lexer.l}, con Flex) y
el analizador sintáctico (\texttt{bison.y}, con Bison) de C-TDS, y dejamos
definida la interfaz del árbol de sintaxis abstracta (AST) que va a sostener
las etapas siguientes del compilador (análisis semántico, generación de
código intermedio y, finalmente, assembly). En esta entrega no implementamos
el AST: dejamos declarada su interfaz en \texttt{ast.h}, con los cuerpos de
\texttt{ast.c} vacíos, y las acciones semánticas de la gramática comentadas,
a la espera de que terminemos de cerrar en equipo algunas decisiones de
diseño que todavía están abiertas.

\section{División del trabajo}

El trabajo de esta etapa se dividió en dos frentes, que avanzaron en paralelo
y se integraron sobre el final:

\begin{itemize}
  \item \textbf{Analizador léxico}: definición y ajuste de las expresiones
  regulares de \texttt{lexer.l} para cubrir todos los tokens de la
  especificación, y de los mensajes de error léxico.
  \item \textbf{Analizador sintáctico y AST}: diseño de la gramática de
  \texttt{bison.y}, detección y resolución del conflicto shift/reduce
  descripto en la Sección~\ref{sec:decisiones}, y diseño de la interfaz del
  AST en \texttt{ast.h}.
\end{itemize}

Las decisiones de diseño más importantes (la resolución del conflicto
shift/reduce y el criterio para el AST) se discutieron y acordaron en
conjunto entre los tres integrantes antes de llevarlas al código, para no
depender de que una sola persona las resolviera por su cuenta.

\section{Analizador léxico}

El lexer reconoce todos los tokens que describe la especificación del
lenguaje: las palabras reservadas (\texttt{boolean}, \texttt{int},
\texttt{float}, \texttt{void}, \texttt{if}, \texttt{else}, \texttt{while},
\texttt{return}), los identificadores, las constantes numéricas (enteras y de
punto flotante), las constantes booleanas, los operadores aritméticos,
relacionales y lógicos, los delimitadores, y los comentarios de línea y
multilínea.

\subsection{\texttt{main} no es palabra reservada}

Decidimos no tratar \texttt{main} como palabra reservada, a pesar de que el
lenguaje exige que todo programa tenga un método con ese nombre. La razón es
que esa exigencia es una regla semántica (debe existir un método llamado
\texttt{main} con una firma particular), no una regla léxica ni sintáctica:
en la gramática, \texttt{main} es simplemente el identificador de un método
como cualquier otro. Si lo hubiéramos declarado como palabra reservada, un
identificador de esa forma dejaría de poder usarse como nombre de variable o
de cualquier otro método, lo cual no está pedido en ningún lado de la
especificación. Por eso \texttt{main} matchea la regla general de
\texttt{ID}, y la verificación de que exista un método \texttt{main} con la
firma correcta queda para el análisis semántico.

\subsection{Número de línea en los mensajes de error}

Activamos \texttt{\%option yylineno} y lo usamos tanto en los mensajes de
error léxico del propio \texttt{lexer.l} como en \texttt{yyerror} del lado de
Bison. Nos pareció imprescindible para que un mensaje de error sea útil: sin
número de línea, encontrar dónde está el problema en un archivo fuente de
varias decenas de líneas obliga a revisar todo el archivo a mano.

\section{Decisiones de diseño, asunciones y extensiones}
\label{sec:decisiones}

\subsection{Un conflicto shift/reduce que nos obligó a rediseñar la gramática}

La primera versión de la gramática que escribimos separaba las declaraciones
globales en dos listas independientes, una para variables y otra para
métodos:

\begin{verbatim}
Program : VariableDeclarations MethodDeclarations ;
\end{verbatim}

Esta estructura calca directamente el orden en que aparecen en la letra los
no terminales \texttt{$\langle$var\_decl$\rangle$} y
\texttt{$\langle$method\_decl$\rangle$}, pero no funciona en la práctica: al
probarla contra el programa de ejemplo de la propia especificación del
lenguaje, que declara una variable global y a continuación un método con
retorno no-\texttt{void},

\begin{verbatim}
int a;
int sumar(int a, int b) { }
\end{verbatim}

el parser fallaba con un error de sintaxis apenas encontraba el segundo
\texttt{int}. La causa es un conflicto de desplazamiento/reducción
(shift/reduce) inherente a un parser LALR(1), que solo tiene un token de
adelanto (lookahead) para decidir. Tanto una declaración de variable como una
de método empiezan igual, con \texttt{Type ID}, y Bison, ante la duda,
siempre prefiere desplazar antes que reducir: al ver el segundo \texttt{int}
todavía dentro de \texttt{VariableDeclarations}, asumía que se trataba de otra
variable y esperaba un \texttt{;} o una \texttt{,}. Como en cambio encontraba
un \texttt{(}, la lista de variables nunca llegaba a cerrarse y el parser no
tenía forma de retroceder para reinterpretar esa declaración como un método.
En la práctica, esto significa que ningún método con tipo de retorno
no-\texttt{void} lograba parsear si estaba precedido por al menos una
variable global, lo cual invalidaba directamente el ejemplo de la letra.

La solución que aplicamos fue unificar ambas listas en una sola, dejando que
la decisión de si \texttt{Type ID} arranca una variable o un método se tome
recién al ver el token siguiente:

\begin{verbatim}
Declarations
    : /* vacio */
    | Declaration Declarations
    ;

Declaration
    : VariableDeclaration
    | MethodDeclaration
    ;
\end{verbatim}

Con esta forma, el parser consume \texttt{Type ID} de manera neutra, sin
comprometerse todavía a ninguna de las dos alternativas, y recién decide al
ver el token siguiente: si es \texttt{;} o \texttt{,}, se trata de una
variable; si es \texttt{(}, de un método. Esto permite además que variables y
métodos aparezcan intercalados en el scope global, que es exactamente lo que
hace el ejemplo de la letra. De paso, aprovechamos para sacar el no terminal
intermedio \texttt{ReturnType} e integrar directamente \texttt{Type} y
\texttt{VOID} en \texttt{MethodDeclaration}, para no introducir una reducción
adicional que pudiera volver a confundir al parser en ese mismo punto.

\subsection{Precedencia de operadores}

Para resolver la precedencia y asociatividad de los operadores de
\texttt{Expression} usamos las declaraciones \texttt{\%left} y
\texttt{\%right} de Bison, siguiendo la tabla de precedencia de menor a mayor
que se desprende de la especificación del lenguaje: disyunción lógica,
conjunción lógica, igualdad, relacionales, suma y resta,
multiplicación/división/resto, y por último la negación lógica y el menos
unario. Optamos por esta vía declarativa en lugar de reescribir
\texttt{Expression} como una jerarquía explícita de no terminales por nivel
de precedencia (una alternativa habitual cuando no se cuenta con
declaraciones de precedencia de la herramienta).

\subsection{Comentarios multilínea sin cerrar}

Asumimos, por el momento, que todo comentario multilínea que se abre en el
código fuente se cierra correctamente. El caso de un comentario multilínea
que nunca se cierra (el archivo termina dentro de un comentario abierto)
queda señalado como algo pendiente de resolver en una iteración posterior del
lexer, y se detalla en la Sección~\ref{sec:problemas}.

\subsection{AST: solo interfaz en esta etapa}

Decidimos, como extensión al alcance mínimo pedido para esta entrega, dejar
completamente definida y documentada la interfaz del AST (\texttt{ast.h}),
aunque no forma parte de lo exigido para el analizador sintáctico en sí. La
justificación de esta decisión se desarrolla en la Sección~\ref{sec:ast}.

\section{Descripción de diseño y justificación}
\label{sec:ast}

\subsection{Alternativas consideradas para el conflicto shift/reduce}

Antes de llegar a la unificación de \texttt{Declarations}/\texttt{Declaration}
descripta en la Sección~\ref{sec:decisiones}, consideramos dos alternativas
que descartamos:

\begin{itemize}
  \item \textbf{Aumentar el lookahead}: Bison genera parsers LALR(1), con un
  solo token de adelanto; no hay forma de pedirle más lookahead sin cambiar
  de herramienta o de clase de gramática, así que esta alternativa no era
  viable con lex/yacc.
  \item \textbf{Forzar el orden en la letra} (todas las variables antes que
  los métodos): en teoría resolvía el conflicto, pero invalidaba el propio
  ejemplo de programa que trae la especificación del lenguaje, que intercala
  una variable y un método. La descartamos porque no es fiel al lenguaje que
  tenemos que reconocer.
\end{itemize}

La alternativa que aplicamos (unificar ambas listas en un único no terminal
y diferir la decisión hasta ver el token siguiente a \texttt{Type ID}) es el
patrón estándar para este tipo de ambigüedad de prefijo compartido en
gramáticas LL/LALR, y es la que finalmente seleccionamos porque no exige
tocar el lenguaje ni la herramienta, solo la forma en que factorizamos la
gramática. La verificamos de dos formas: generando la gramática con
\texttt{bison -d -v} y confirmando que el archivo de reporte no muestra
ningún conflicto shift/reduce ni reglas inalcanzables (algo que sí aparecía
con el diseño original), y compilando el parser completo contra Flex y la
interfaz de AST para correrlo sobre el ejemplo de la letra y sobre variantes
propias (variables intercaladas con métodos, condicionales anidados, ciclos,
llamadas a método con y sin argumentos), confirmando que todos los casos
válidos parsean sin errores formales.

\subsection{Por qué separamos el AST de \texttt{bison.y}}

Decidimos no definir la estructura del AST directamente en
\texttt{bison.y}, sino en un archivo aparte, \texttt{ast.h}. La razón es que
sobre este árbol todavía se va a construir el análisis semántico, un
intérprete o generador de código intermedio, y finalmente el generador de
assembly, y no queríamos que esas etapas dependieran de un archivo pensado en
primer lugar para la gramática. Separar la interfaz del árbol en su propio
header nos permite mantener \texttt{bison.y} enfocado en las reglas de
producción y dejar la definición de los nodos como un contrato aparte, que
cualquier etapa posterior puede incluir sin arrastrar nada de Bison.

\subsection{Por qué \texttt{ast.h} solo tiene interfaz, y \texttt{ast.c} está vacío}

Dentro de esa separación, distinguimos además interfaz de implementación:
\texttt{ast.h} declara los tipos y las firmas de las funciones constructoras
del árbol, documentando en comentarios qué hace cada una, quién es dueño de
la memoria que devuelve y qué se espera de sus parámetros; \texttt{ast.c}
contiene esas mismas funciones pero con el cuerpo vacío (un \texttt{TODO}).
No las implementamos todavía porque hacerlo implica tomar decisiones de
diseño que como equipo no terminamos de cerrar (por ejemplo, si el árbol va a
liberar memoria a medida que se descarta, o cómo van a convivir exactamente
los punteros \texttt{left}/\texttt{mid}/\texttt{right} con el arreglo
\texttt{children[]} en los distintos tipos de nodo). Preferimos dejar esas
firmas fijas y bien documentadas ahora, y resolver el cuerpo de cada función
más adelante, cuando todo el equipo esté de acuerdo en cómo hacerlo, en lugar
de que esa decisión quede tomada de forma implícita por quien haya escrito el
primer código. Por eso en esta entrega \texttt{ast.h} es, en los hechos, toda
la documentación de diseño del AST que tenemos: cada función está comentada
con su propósito, sus parámetros y quién es responsable de liberar la memoria
que devuelve, aunque todavía no haga nada.

Por el mismo motivo dejamos las acciones semánticas de \texttt{bison.y}
escritas pero comentadas: cada regla de la gramática indica, a modo de
comentario, cómo se vería la construcción del nodo correspondiente una vez
que el AST esté implementado (qué \texttt{NodeType} le corresponde, qué va en
\texttt{left}/\texttt{mid}/\texttt{right} o en \texttt{children[]}), pero no
se ejecuta ningún código todavía. Esto nos permitió dejar registrada, regla
por regla, la intención de diseño de cada construcción del lenguaje sin
depender de que \texttt{ast.c} ya esté terminado, y sin arriesgarnos a que el
parser genere nodos a medio definir o haga referencias a memoria sin
inicializar. Cuando decidamos avanzar con la implementación del AST, el
trabajo va a consistir en completar los cuerpos de \texttt{ast.c} siguiendo
lo que ya documentamos en \texttt{ast.h}, y recién ahí descomentar estas
acciones.

\subsection{Estructura de \texttt{NodeAST}}

Cada nodo del árbol (\texttt{NodeAST}) tiene un \texttt{NodeType} que indica
qué construcción del lenguaje representa, un \texttt{Symbol} opcional para
guardar el identificador o el valor asociado, punteros \texttt{left},
\texttt{mid} y \texttt{right} para los casos con una cantidad fija y chica de
hijos (por ejemplo, una asignación o un \texttt{if-else}), y un arreglo
\texttt{children[]} de tamaño variable para los casos donde la cantidad de
hijos depende del programa fuente (por ejemplo, la lista de identificadores
de una declaración de variables, o los argumentos de una llamada a método).
El campo \texttt{type} guarda el tipo semántico del nodo (\texttt{TYPE\_INT},
\texttt{TYPE\_BOOL}, etc.), separado de \texttt{nodeType} para no mezclar
``qué construcción sintáctica es este nodo'' con ``qué tipo de dato tiene'', que
son dos preguntas distintas y que además se resuelven en momentos distintos
del proceso (la primera durante el parsing, la segunda recién en el análisis
semántico).

\subsection{Criterio para crear un nodo}

No decidimos crear un nodo por cada no terminal de la gramática, sino
únicamente para las construcciones que aportan información propia o que
cambian el comportamiento del árbol resultante. Hay no terminales que
existen en la gramática solo para organizar la sintaxis (evitar ambigüedades,
factorizar alternativas) pero que en el AST no necesitan representación
propia: por ejemplo, \texttt{Statement} agrupa todas las variantes de
sentencia posibles, pero el nodo real lo termina aportando la alternativa
concreta que matcheó (una asignación, un \texttt{return}, un bloque), así que
\texttt{Statement} simplemente redirige el valor de su hijo hacia
\texttt{\$\$} sin crear un nodo nuevo. En cambio, \texttt{Block} sí tiene su
propio \texttt{BLOCK\_NODE}, porque no es un simple contenedor sintáctico:
representa la apertura de un nuevo scope y agrupa las declaraciones de
variables y las sentencias que contiene, que sí son información real que el
árbol necesita conservar.

\section{Detalles de implementación interesantes}

\subsection{Resolución de la ambigüedad \texttt{Type ID}}

El punto más particular de esta etapa es, en términos de implementación, la
técnica usada para resolver la ambigüedad de prefijo compartido entre
\texttt{VariableDeclaration} y \texttt{MethodDeclaration} (ver
Sección~\ref{sec:decisiones}). No es un algoritmo en el sentido clásico, sino
una técnica de factorización de gramáticas: en vez de que dos producciones
alternativas empiecen ambas por los mismos símbolos y dejen que el parser
``adivine'' cuál es, se las fusiona en una sola producción que consume el
prefijo común y recién más adelante, con la gramática ya en un estado neutro,
se bifurca según el símbolo siguiente. Es el mismo principio que la
factorización por la izquierda en gramáticas LL, aplicado acá para evitarle
al parser LALR(1) una decisión que no puede tomar con un solo token de
adelanto.

\subsection{Listas de hijos mediante \texttt{NodeList}}

Para los nodos con una cantidad variable de hijos (por ejemplo,
\texttt{VariableDeclaration}, cuya cantidad de identificadores declarados no
se conoce de antemano), la interfaz de \texttt{ast.h} separa dos
responsabilidades: \texttt{NodeList} es una lista enlazada transitoria, que
solo existe durante la construcción del árbol dentro de \texttt{bison.y} para
ir acumulando nodos hermanos a medida que la gramática reduce reglas
recursivas (\texttt{IdentifierList}, \texttt{VariableDeclarations},
\texttt{Statements}); y \texttt{attachChildren} es la función encargada de
volcar esa lista al arreglo \texttt{children[]} del nodo padre una única vez,
dejando el AST final compuesto solo por \texttt{NodeAST}, sin rastros de la
lista auxiliar. Elegimos este esquema en dos pasos (acumular en lista,
después volcar a arreglo) en lugar de, por ejemplo, usar directamente un
arreglo dinámico que se va agrandando durante el parsing, porque encaja mejor
con la forma en que Bison arma los valores semánticos de abajo hacia arriba
en reglas recursivas por la derecha.

\section{Problemas conocidos}
\label{sec:problemas}

\subsection{Comentario multilínea sin cerrar}

\textbf{Detección}: al construir casos de test para el lexer, probamos
deliberadamente un archivo fuente donde un comentario multilínea se abre con
\texttt{/*} y el archivo termina sin el \texttt{*/} correspondiente. La
expresión regular actual de \texttt{lexer.l} para comentarios multilínea
exige que el cierre exista, así que en este caso no matchea ninguna regla de
comentario.

\textbf{Caso de test}:
\begin{verbatim}
/* este comentario
nunca se cierra
int nuncaSeEjecuta;
\end{verbatim}

\textbf{Comportamiento actual}: como la regla de comentario no matchea (no
hay \texttt{*/}), el lexer cae en las reglas siguientes y termina
interpretando el contenido del comentario carácter por carácter, reportando
cada símbolo dentro de él como error léxico, en lugar de señalar el problema
real, que es que el comentario nunca se cerró.

\textbf{Acción tomada}: identificamos el problema y la solución conceptual
(introducir un estado exclusivo de Flex, por ejemplo \texttt{\%x COMMENT},
que se entre al ver \texttt{/*} y se salga al ver \texttt{*/}, junto con una
regla de \texttt{<<EOF>>} dentro de ese estado que reporte explícitamente
``comentario no cerrado''), pero decidimos no aplicarla todavía en
\texttt{lexer.l}. Queda pendiente para una iteración posterior del lexer.

\subsection{Alcance de esta entrega}

Del AST, en esta etapa solo dejamos definida la interfaz (\texttt{ast.h}) y
las invocaciones ya escritas pero comentadas dentro de \texttt{bison.y}, a
modo de referencia de cómo se va a construir cada nodo; los cuerpos de
\texttt{ast.c} están vacíos. No es un problema en el sentido de algo que
falle: es una decisión deliberada (ver Sección~\ref{sec:ast}), porque
implementar el AST implica decisiones de diseño que todavía tenemos que
cerrar como equipo. La próxima etapa consiste en completar los cuerpos de
\texttt{ast.c} y descomentar las acciones semánticas de \texttt{bison.y} una
vez que terminemos de acordar esos últimos detalles.

\end{document}
